% Lliçó: àlgebra de límits (valencià).

\begin{frame}
\didactatitle{Àlgebra de límits}

\begin{proposition}[Àlgebra de límits]
Si $\lim_{x \to a} f(x) = L$ i $\lim_{x \to a} g(x) = M$, aleshores
\begin{itemize}
  \item $\lim_{x \to a} \bigl(f(x) + g(x)\bigr) = L + M$;
  \item $\lim_{x \to a} f(x)\,g(x) = L M$;
  \item $\lim_{x \to a} \dfrac{f(x)}{g(x)} = \dfrac{L}{M}$, si $M \neq 0$.
\end{itemize}
\end{proposition}

\begin{notesonly}
\begin{proof}
Provem la primera; les altres dues segueixen la mateixa idea amb alguna
fitació més. Donat $\varepsilon > 0$, existeixen $\delta_1, \delta_2 > 0$
tals que $|f(x) - L| < \varepsilon/2$ si $0 < |x - a| < \delta_1$, i
$|g(x) - M| < \varepsilon/2$ si $0 < |x - a| < \delta_2$. Prenent
$\delta = \min\{\delta_1, \delta_2\}$, per a $0 < |x - a| < \delta$ la
desigualtat triangular dona
\[
  \bigl|\bigl(f(x) + g(x)\bigr) - (L + M)\bigr|
  \le |f(x) - L| + |g(x) - M|
  < \tfrac{\varepsilon}{2} + \tfrac{\varepsilon}{2} = \varepsilon .
\]
\end{proof}
\end{notesonly}

\dpause

\begin{remark}
Com que $\lim_{x \to a} x = a$ i el límit d'una constant és ella mateixa, la
proposició dona el límit d'un polinomi avaluant: $\lim_{x \to a} p(x) =
p(a)$. I el d'una funció racional, $\lim_{x \to a} p(x)/q(x) = p(a)/q(a)$,
sempre que $q(a) \neq 0$.
\end{remark}

\timing{15 min}
\end{frame}
